\section{结果与讨论}\label{sec:results}
\subsection{示例数据}
% !lp begin {"id":"results-table","title":"阅读记录汇总","status":"draft","tag":"人工初稿","note":"演示数据，可尝试编辑表格"}
表\ref{tab:reading}给出虚构的一周记录。工作日与周末仅用于展示分类，数据不能用于实际研究结论。
\begin{table}[htbp]
\centering
\caption{一周阅读记录示例（分钟）}\label{tab:reading}
\begin{tabular}{lrrr}
\toprule
时间 & 天数 & 平均时长 & 总时长 \\
\midrule
工作日 & 5 & 36 & 180 \\
周末 & 2 & 60 & 120 \\
合计 & 7 & 42.9 & 300 \\
\bottomrule
\end{tabular}
\end{table}
% !lp end
\subsection{结果解读}
% !lp begin {"id":"results-discussion","title":"结果解读","status":"draft","tag":"AI初稿","note":"检查解释是否有证据支撑"}
在这组虚构数据中，周末的平均阅读时间较长。但由于观察周期很短，我们不能据此推断长期趋势。讨论应区分观测结果、可能解释与尚待验证的假设。
% !lp end
\subsection{局限性}
% !lp begin {"id":"results-limitations","title":"局限性","status":"todo","tag":"待编写","note":"补充样本量和记录偏差"}
此处可以继续补充样本量、记录方式和选择偏差对结论的影响。
% !lp end
